\documentclass[script,con]{debate}
\motion{超级英雄的出现，对世界而言更是灾难（反方）}
\meta{赛制}{中国科学技术大学新生辩论赛（四对四）}
\meta{口径来源}{备赛文档 第十节 10.2 反方口径表。本文稿所有定义、判准、论点、数据均取自该表}
\begin{document}
\debatetitle
\begin{thesisbox}[证据纪律]漫威、DC 的情节是设定，不是论据。全场只用现实世界的对应物。本方最容易犯的错就是拿电影里的反派论证风险，一旦被正方点出“这是编剧写的”，整条证据链都会显得不严肃\end{thesisbox}
\begin{speech}{1. 反一开篇立论稿}{正文 839 字 / 180 秒}
谢谢主席，各位评委。

刚才的质询里，正方一辩已经承认，超级英雄不由现有法律和权力机构授权，也不由它们问责。请评委记住这句话，我方全场都建立在它上面。

定义上我方和对方没有分歧。超级英雄，指拥有远超常人的能力、主动介入公共安全事务的人，他行动的依据是自己的判断，不受既有权力的授权与问责。出现，指这种人存在于世界上成为事实，并且被世界知道。至于灾难，我方承认它不等于“有风险”，它指的是世界的处境重大地变糟，而且改不回来。

判准我方也接受对方的版本：长期看普通人的处境是变好还是变坏，看安全和自主两项。分歧在于对方只算了救人这一栏。我方要证明的是，另一栏的损失是结构性的、改不回来的，而且大到足以把救人那一栏抵消。

我方第一点，没人能拦。

韦伯一九一九年讲过一句话：使用物理暴力的权利，只在国家允许的限度之内被赋予其他机构或个人。现代世界的安全底座就是这句话。超级英雄是这条线外的第一个例外。有人会说这只是理论。那我举个例子。二〇一一年十月的西雅图，一个真的穿上制服上街打击犯罪的年轻人，把一群刚从夜店出来的人当成了斗殴，朝他们喷了胡椒喷雾。西雅图警方说，他打扮成什么样都不代表他可以凌驾于法律之上。而他连超能力都没有。

我方第二点，坏人也会有。

能力什么样，和拿着它的人什么样，本来就是两件事。一旦世界知道这种能力可以落在个体身上，各国的第一反应不是欢呼，是找到它，或者抢在别人之前先下手。人类上一次遇到压倒性的力量，八十年过去，留下一万两千多枚核弹头。你想想一九八三年九月的那个晚上，苏联的预警系统报告美国打来五枚导弹，值班军官凭一句“真开战不会只打五枚”没有上报。后来查明，是阳光照在高空云层上的反射。人类和那一夜之间隔着的，是一个人的判断。

我方第三点，等靠退化。

有人兜底，人就会松手，这不是懦弱，是理性。二〇〇五年有一项实验，一百一十二个学生，被要求写“超人”的那一组，当场报名做志愿者的只有百分之二十三；写“超级英雄”这个类别的有百分之四十二。九十天后真正到会场的，前者百分之四，后者百分之十七。制度也一样。全球正在威权化的四十五个国家里，二十七个起点是民主国家，到二〇二四年只剩九个还是。人先松手，再发现没人接着。

所以我方全场只请对方回答一个问题：他错了以后，谁能把它改回来？谢谢。

\begin{contingency}
\ifthen{正一立论没说“不受问责”，就在反二驳论第一句点出他藏了这一条。}{}
\ifthen{正方把判准改成“只看救了多少人”，就用“那是半本账”一句话带回。}{}
\end{contingency}
\end{speech}
\begin{stage}{2. 反一被质询防守预案}{正四质询，120 秒双边计时}
回答原则：短，不否认常识，不滑到“他一定会作恶”。被逼二选一时指出预设有问题，再给我方的第三种说法。

\begin{qa}[对方会问][我方回答（简短，守住前提）]
\qarow{你说的灾难，是和哪个世界比？}{和他出现之前的世界比。之前世界有纠错能力，之后没有了。}
\qarow{你是不是在和一个完美制度比？}{不是。我方比的是纠错能力有没有，不是它好不好用。}
\qarow{那卢旺达那一百天，谁纠正了？}{没有人。我方承认那是失败。失败可以追责，可以改制度，这就是纠错。}
\qarow{你是不是主张他一定会作恶？}{不主张。我方的问题是无法验证，也无法在他错的时候改回来。}
\qarow{那你只证明了风险，不是灾难？}{纠错能力的消失不是风险，是已经发生的事实。}
\qarow{消防员在火场里谁能当场问责他？}{没人。但他事后能被开除、被起诉、被换掉。超级英雄不能。}
\qarow{你举的义警都没有超能力，怎么类比？}{我方类比的是能力和问责之间的落差，不是能力大小。}
\qarow{落差更大就更坏吗？}{落差更大，出错的代价更大，能纠正的可能更小。}
\qarow{你引用的那项研究没有对照组，你知道吗？}{知道，作者自己写明了。所以我方的主证据是制度那一组数据。}
\qarow{你承不承认有些人本来会死？}{承认。我方从不否认救援的价值。}
\end{qa}
\end{stage}
\begin{stage}{3. 反四质询正一问题链}{120 秒，双边计时}
\begin{chain}{链一：把“不受问责”钉成正方自己的承认}{拿到全场最关键的一句承认}
\q{超级英雄受不受现行法律的问责？}{答“不受”：感谢，这句话我方会用一整场。 → 答“受”：那他就是执法者，请问归哪个机关管？}
\q{他决定今天救谁、不救谁，要不要经过任何人同意？}{答“不要”：那这个决定权属于谁？}
\q{如果他判断错了，把无辜的人打伤，谁能强制他停下来？}{答“舆论”：舆论对一个不需要任何人配合的人有强制力吗？}
\cut{好的，我理解为您没有否认。下一个问题。}
\close{感谢。正方已经承认，没有任何人能强制他停下来。}
\end{chain}
\begin{chain}{链二：比较对象}{拆掉正方“跟现实比”的框，暴露他们只算半本账}
\q{对方的判准里有“自主”这一项，对吗？}{承认：那纠错能力算不算自主的一部分？}
\q{对方今天有没有给纠错能力这一栏记过任何代价？}{答“记了”：请具体说数字。 → 答“没有”：那是半本账。}
\q{一百天里的一百万人，是缺能力还是缺授权？}{}
\close{对方自己承认，卢旺达缺的是决心。}
\end{chain}
\begin{chain}{链三：证据来源}{堵死正方用设定当论据，同时约束我方自己}
\q{对方今天会不会用电影情节论证现实结论？}{答“不会”：很好，双方共同遵守。}
\q{对方有没有一项研究，证明不可问责的强者会长期做好事？}{答“没有”：那双方在这一点上都只有类比。}
\close{对方承认了，他们那一栏没有直接证据。}
\end{chain}
\end{stage}
\begin{speech}{4. 反二驳论稿}{正文 557 字 / 120 秒}
正方立论有一个共同的毛病：他们把账算了一半。

先说正方第一点，救得下来。他们举了切尔诺贝利，举了卢旺达。可请评委回想一下卢旺达的细节：联合国维和部队从两千一百多人减到两百七十人，是因为会员国不肯增派部队、不肯加强授权。这是联合国自己的说法。那一百天缺的不是能力，是有人肯不肯下那个决定。切尔诺贝利也一样，当时的机械在高辐射下失灵，这是技术问题，而技术后来补上了。正方把“没人肯做”和“没人能做”混成了一件事，然后说超级英雄能补。可他补上的，只是把“没人决定”换成“一个人决定”。

再说正方第二点，恶做不成。正方引用的威慑研究，说被抓住的确定性比惩罚更有威慑力；正方引用的那六十五项研究，讲的是把警力放到犯罪高发的小片区域。请评委注意，这两处研究里的人，全都是穿制服、有编号、能被投诉和起诉的警察。正方把它们外推到一个不受问责的个体身上，这一步没有论证。更要紧的是，同样的逻辑用在他身上，“被抓住的确定性”就变成了“他今天心情如何的确定性”。威慑越强，这件事越危险，因为方向由他一个人定。

至于正方第三点，正方一辩自己说，同一篇论文里还有一个相反的结果。是的，那就是我方立论的数据：写“超人”的那一组，九十天后到场率百分之四，写类别的那一组百分之十七。我方承认那项研究没有设中性对照组，作者写明了。所以我方把它当成机制的说明，我方的主证据是四十五个威权化国家里，二十七个起点是民主国家，只剩九个还是。

正方的问题，从头到尾是同一个：他们不给纠错能力这一栏记任何代价。

\begin{contingency}
\ifthen{正二先回应我方而不立新点：把时间用在追问“纠错能力那一栏你记了多少”。}{}
\ifthen{正方拿出新数据：先问出处与年份，再问它研究的对象是不是可被问责的人。}{}
\end{contingency}
\end{speech}
\begin{stage}{5. 反二对辩要点}{90 秒，单边计时}
进攻点（每个 15 到 20 秒，说完就停）：

\begin{enumerate}[leftmargin=1.8em,itemsep=2pt,topsep=2pt]
\item 纠错能力这一栏，你方记了多少代价？请给一个说法。
\item 你引用的威慑研究，研究对象是不是可以被投诉和起诉的警察？
\item 卢旺达缺的是能力还是决心？联合国自己怎么说的？
\item 他错了以后，谁能强制他停下来？这是我方第二次问。
\item 你方一辩承认他不受问责，那这个前提对世界意味着什么，请正面说。
\end{enumerate}
备用点（用于回应对方进攻）：

\begin{enumerate}[leftmargin=1.8em,itemsep=2pt,topsep=2pt]
\item 我方不主张他会作恶。我方主张我们无法验证，也无法改回来。
\item 义警的类比落点是问责落差，不是能力大小。能力越大，落差越大。
\item 消防员事后能被开除、被起诉、被换掉，这就是区别。
\item 那项研究没有对照组，我方自己先说的，不是被你们抓到的。
\item 一万两千多枚核弹头不是“人类学会共存”的证据，是人类反应的证据。
\item 一九八三年九月那一夜，救下世界的是一个人违反程序，这恰恰吓人。
\end{enumerate}
对方最可能问的三个问题：

\begin{bullets}
\item “你只证明了风险，不是灾难？” → 纠错能力的消失是已经发生的事实，不是概率。
\item “你凭什么假设他是坏人？” → 我方从不假设。我方的问题是无法验证。
\item “那些本来会死的人怎么办？” → 我方承认他们的命是真的。请你也承认另一栏的代价是真的。
\end{bullets}
\end{stage}
\begin{stage}{6. 反三盘问问题链}{90 秒，单边计时}
\begin{chain}{链一：问二辩再问四辩，检验口径}{抓正方对“纠错”的两套说法}
\q{问正二：你方是否承认，他错了以后没有人能强制他停下来？}{}
\q{同一个问题问正四：你方是否承认这一点？}{}
\q{如果两人答案不同：究竟有没有人能强制他停下来？}{}
\q{如果两人都承认：对方给“自主”这一栏记过代价吗？}{}
\end{chain}
\begin{chain}{链二：把救援与授权分开}{为小结准备一句话}
\q{联合国在卢旺达的记载里，缺的是部队还是会员国的授权？}{}
\q{如果缺的是授权，那超级英雄解决的是哪一个问题？}{}
\q{对方有没有一个案例，是超常能力本身把一场灾难挡下来了？}{}
\end{chain}
盘问目标（小结里要证明什么）：正方全场没有给纠错能力记任何代价；正方最好的两个案例，缺的都是决心与授权，不是能力。

\end{stage}
\begin{stage}{7. 反二、反四被盘问防守预案}{两人口径必须一致}
\begin{qaa}[对方会问][二辩答][四辩答（必须一致）]
\qaarow{你方是否主张他迟早会作恶？}{不主张。我方主张无法验证，也无法纠错。}{不主张。同样一句话：无法验证，无法纠错。}
\qaarow{那你只证明了风险？}{纠错能力的消失是已经发生的事实。}{是事实，不是概率。}
\qaarow{你说的灾难是和哪个世界比？}{和他出现之前的世界比。}{和他出现之前的世界比。}
\qaarow{那个世界里卢旺达谁纠正了？}{没有人。我方承认那是失败，但失败能追责、能改制度。}{没有人。可那次失败之后，联合国做了独立调查。}
\qaarow{义警没有超能力，怎么类比？}{类比的是问责落差，不是能力大小。}{同上。落差是共同点。}
\qaarow{你引的研究没有对照组？}{是，作者写明了，我方自己先说的。}{是。所以主证据是制度那一组数据。}
\qaarow{你承不承认有人本来会死？}{承认。我方从不否认救援的价值。}{承认。我方请对方也承认另一栏。}
\end{qaa}
\end{stage}
\begin{speech}{8. 反三盘问小结稿}{正文 327 字 + 临场位 80 字 / 90 秒}
\reserve{此处留临场位约 80 字：直接回应正三小结里最新出现的说法}
谢谢主席。刚才的盘问我只想请评委看两件事。

第一件，我问正方二辩和四辩同一个问题：他错了以后，有没有人能强制他停下来。两个人都答不出一个能执行的办法。那么对方判准里写着的“自主”这一项，今天到底记了多少代价？到现在为止是零。这就是我方说的半本账。

第二件，我请对方举一个案例，是超常能力本身把一场灾难挡下来了。对方举不出来。对方举的是卢旺达，可联合国自己的记载写得清清楚楚：兵力从两千一百多人减到两百七十人，原因是会员国不肯增派、不肯加强授权。老实说，那一百天里最缺的东西，不是力气。

各位评委，这一栏不是我方刁难对方，它写在双方都接受的判准里。我方承认，对方救下来的那些人是真的，可一本账要两栏一起算。

所以现在的分歧很清楚了。对方在说，多一个能救人的人总比没有好。我方在说，请你把纠错能力那一栏也填上。空着不填，那不叫赢，叫没算完。谢谢。

\begin{contingency}
\ifthen{正三小结已经先回答了“谁能强制他停下来”，就在临场位里直接拆那个答案，不要重复问题。}{}
\end{contingency}
\end{speech}
\begin{stage}{9. 自由辩战场设计}{240 秒}
\begin{battleground}{战场一：谁能把它改回来}
\begin{bullets}
\item 目标：让评委听见正方第四次给不出一个能执行的办法。
\item 开场问题：他错了以后，谁能强制他停下来？
\item 对方答“舆论、道德” → 追问：舆论对一个不需要任何人配合的人有强制力吗？
\item 对方答“没有人，但净变化是正的” → 追问：那纠错能力这一栏你记了多少代价？
\item 对方可能反攻：消防员也不能当场被问责 → 我方防守：消防员事后能被开除、被起诉、被换掉，这就是区别。
\item 收束语：一整场比赛，对方没有给出一个能执行的办法。
\end{bullets}
\end{battleground}
\begin{battleground}{战场二：缺的是能力还是决心}
\begin{bullets}
\item 目标：把正方最好的两个案例从“能力缺口”改判成“授权缺口”。
\item 开场问题：联合国自己的记载里，卢旺达缺的是部队还是授权？
\item 对方答“两者都缺” → 追问：那超级英雄补的是哪一个？
\item 对方答“缺能力” → 追问：请对方念一下联合国那句话，会员国不肯做什么？
\item 对方可能反攻：切尔诺贝利只能用人去顶 → 我方防守：那是技术在高辐射下失灵，技术后来补上了，而不可问责补不上。
\item 收束语：对方最有力的两个案例，缺的都不是力气。
\end{bullets}
\end{battleground}
\begin{battleground}{战场三：外推是否成立}
\begin{bullets}
\item 目标：拆掉正方威慑论据的适用范围。
\item 开场问题：你引用的威慑研究里，研究对象是不是可以被投诉和起诉的警察？
\item 对方答“是” → 追问：那把结论搬到一个不可问责的人身上，这一步谁论证？
\item 对方答“机制一样” → 追问：机制一样的话，威慑的方向由谁决定？
\item 收束语：同一套逻辑放到他身上，确定性就变成了他一个人的确定性。
\end{bullets}
必问（前两分钟问完）：

\begin{bullets}
\item 他错了以后，谁能强制他停下来？
\item 纠错能力这一栏你方记了多少代价？
\item 卢旺达缺的是能力还是授权？
\item 你引用的威慑研究，研究对象可不可以被起诉？
\end{bullets}
必答与口径：

\begin{bullets}
\item “你只证明了风险？” → 纠错能力的消失是已经发生的事实，不是概率。
\item “你凭什么假设他是坏人？” → 我方从不假设，我方的问题是无法验证。
\item “那些本来会死的人呢？” → 承认他们的命是真的，也请你承认另一栏是真的。
\end{bullets}
时间分配与站位：前 90 秒打战场一并把必问抛完；中间 90 秒打战场二；最后 60 秒打战场三并收束，不开新战场，留 20 秒备用。一辩守定义与判准，二辩接驳论过的点，三辩接盘问成果，四辩负责收束与价值。

\end{battleground}
\end{stage}
\begin{speech}{10. 反四结辩稿}{正文 899 字 / 210 秒，封路式}
谢谢主席，各位评委。

我先说一件正方四辩等一下很可能会做的事。他会告诉各位，我方在拿一个不存在的完美制度和现实比；他会说风险是关于将来的，而那些人的死是正在发生的。这条路我现在就走一遍，因为等他走完，我没有机会再说话了。

我方从来没有拿完美制度作比较。我方比的是他出现之前和之后。之前的世界会失败，卢旺达那一百天就是失败。可失败之后，联合国做了独立调查，会员国被点名，规则被改写。那叫纠错。而在超级英雄出现之后，他做错的那一次，没有调查，也不会有下一次改进。请评委注意，我方要的从来不是一个不会出错的世界，是一个出错之后能改回来的世界。

正方全场最有力的是那句“有些伤害只差那几分钟”。我承认这句话打动人。我也承认，没有超级英雄的世界里确实有人救不下来。可正方自己举的两个例子，替我方说了话。卢旺达那一百天，联合国的记载是兵力从两千一百多人减到两百七十人，原因是会员国不肯增派、不肯加强授权。缺的是决心。切尔诺贝利现场只能用人的身体去顶，是因为当时的机械在高辐射下失灵，而机械后来补上了。缺的是技术。正方的那几分钟，从来不是没有力气，是没有人下决定，或者机器还没造出来。超级英雄换掉的，只是“谁来下这个决定”。

那就回到我方一辩问了一整场的问题：他错了以后，谁能把它改回来？正方二辩说不出来，四辩也说不出来。到现在为止，正方判准里“自主”那一栏还是空的。

我方的三点还立着。韦伯说过，使用物理暴力的权利只在国家允许的限度内被赋予个人，这是现代世界的安全底座，而他是底座之外的第一个例外。人类上一次面对压倒性的力量，八十年下来留下一万两千多枚核弹头，一九八三年九月那一夜，隔在人类和核战争中间的是一个军官不肯按程序办事。至于第三点，四十五个正在威权化的国家里，二十七个起点是民主国家，只剩九个还是。人先松手，再发现没人接着。

最后我想说一件小事。

二〇一一年十月一个凌晨，西雅图一家夜店门口，几个人刚出来，正在跳舞。一个穿着制服、自称在打击犯罪的年轻人从后面走过来，觉得他们在打架，朝他们脸上喷了胡椒喷雾。那几个人蹲在地上，睁不开眼。他们没有做错任何事。

我一直在想，那天晚上他们心里最先冒出来的念头是什么。大概不是疼。是“怎么会这样”，和“现在该找谁”。

后来警方没能起诉他，因为找不齐当事人。

那个年轻人连超能力都没有。

我方今天争的不是世上该不该有好人。是当一个人错了，另一个人还能不能找到地方说理。请各位把票投给反方。谢谢。

\begin{contingency}
\ifthen{正四把战场拉到“你只证明了风险”：用“纠错能力的消失是已经发生的事实”一句话回。}{}
\ifthen{正方在自由辩已经答出了一个纠错办法：把结辩第四段改成拆那个办法，其余不动。}{}
\end{contingency}
\end{speech}
\end{document}
